\bookchapter{Capítulo 44: Prática do 0 ao 1 — o fluxo completo de um projeto inteiro}

\chapterintro{A teoria acabou; agora é hora de colocar a mão na massa. Este capítulo não é um tutorial simples de 'digitar junto', mas uma metodologia prática de 'como pensar, como pedir e como validar'. Usaremos três projetos de dificuldades diferentes para mostrar o processo completo de raciocínio, dos requisitos ao produto final. Não importa qual ferramenta você use: esse fluxo serve para todos.}

\booksection{Projeto 1: site pessoal minimalista de ferramentas (nível iniciante)}

\booksubsection{Pensando nos requisitos}

Objetivo: criar uma página pessoal de ferramentas que exiba os atalhos para 3 ou 4 utilitários práticos. É um projeto puramente de front-end, sem back-end nem banco de dados. O ponto principal não é a complexidade das ferramentas, mas vivenciar o fluxo completo 'dos requisitos ao produto final'.

Seguindo os cinco elementos de requisitos: Who (quem) — uso pessoal. What (o quê) — uma página que exibe os atalhos das ferramentas. Why (por quê) — reunir as ferramentas mais usadas em uma única página para facilitar o acesso. How (como) — um único arquivo HTML+CSS+JS, sem frameworks. Constraint (restrição) — pronto em até 30 minutos, sem precisar fazer deploy.

\booksubsection{Passando as instruções}

Seguindo o princípio Think Before Coding, a sua instrução deve conter informações suficientes para a IA entender o objetivo, sem precisar especificar cada detalhe de implementação:

\begin{lstlisting}
Crie uma página web de caixa de ferramentas pessoal, requisitos:
1. HTML em arquivo único (CSS e JS inline), sem frameworks
2. Tema escuro, estilo minimalista moderno
3. Mostrar 4 cards de ferramentas:
   - Resumidor de texto
   - Contador de palavras
   - Seletor de cores
   - Gerador de QR code
4. Cada card tem ícone, título e descrição curta
5. Os cards têm animação ao passar o mouse
6. Layout responsivo, utilizável no celular
7. Uma barra de título no topo
Dê diretamente o arquivo HTML completo.
\end{lstlisting}

\booksubsection{Aceitação e iteração}

Depois que a IA entregar o código, faça a aceitação com as três camadas de verificação da iteração. Primeira camada, verificação funcional: o arquivo abre no navegador? Os 4 cartões aparecem? A animação ao passar o mouse funciona? Segunda camada, verificação de limites: o layout fica correto no tamanho de tela de celular? Os cartões quebram a linha quando a janela é muito estreita? Terceira camada, verificação de arquitetura: o código é legível? O CSS está embutido em uma única tag style? Há redundâncias evidentes?

Se você também quiser implementar a função de algum cartão de ferramenta (e não apenas o atalho), adicione aos poucos — implemente a função de uma ferramenta por vez, valide e só depois adicione a próxima. É a aplicação prática da estratégia de 'passos pequenos e rápidos'.

\booksection{Projeto 2: sistema de controle financeiro full stack (nível intermediário)}

\booksubsection{Requisitos e PRD}

Objetivo: uma aplicação de controle financeiro full stack, com front-end, back-end e banco de dados. Este projeto é bem mais complexo que o primeiro e vai usar os métodos de engenharia de requisitos do Capítulo 40.

Comece escrevendo o PRD. Usando o modelo da Seção 3 do Capítulo 40, escreva a visão geral do projeto, a lista de funcionalidades, a solução técnica, o modelo de dados e o plano de desenvolvimento. Na escolha da solução técnica, considere o princípio Simplicity First — front-end em React, mas sem gerenciamento de estado complexo como o Redux; back-end em Node.js, mas sem microsserviços; banco de dados SQLite em vez de PostgreSQL (na fase de desenvolvimento).

O plano de desenvolvimento se divide em 5 passos: Passo 1, esqueleto da API no back-end (servidor Express + SQLite + endpoints CRUD); Passo 2, páginas básicas do front-end (lista de lançamentos + formulário de adição); Passo 3, integração entre front-end e back-end (chamadas de API + exibição dos dados); Passo 4, categorias e filtros; Passo 5, resumo mensal e gráficos simples. Ao fim de cada passo, você tem algo 'que funciona'.

\booksubsection{Criando o arquivo de memória do projeto}

Antes de começar a programar, crie o arquivo de memória do projeto. Essa é a prática central da engenharia de contexto — fazer a IA conhecer o histórico do projeto no início de cada conversa.

\begin{lstlisting}
# Projeto: app de finanças full-stack
## Visão geral
App de finanças pessoais com registros de receitas/despesas, categorias e totais mensais.
## Stack tecnológica
- Frontend: React + TypeScript + Vite
- Backend: Node.js + Express + TypeScript
- Banco de dados: SQLite (via better-sqlite3)
- Definições de tipo compartilhadas entre frontend e backend
## Convenções de API
- Caminho base: /api
- Formato de resposta:{ success: boolean, data: any, error: string }
- Estilo RESTful
## Modelo de dados
Record { id, amount, type, category, note, date }
## Padrões de código
- Modo estrito do TypeScript
- Indentação de 2 espaços
- Funções com comentários JSDoc
- Comentários e mensagens de commit em português
\end{lstlisting}

\booksubsection{Implementando passo a passo}

Siga o plano de desenvolvimento usando uma conversa independente para cada passo. Abaixo, um exemplo de instrução para o primeiro passo:

\begin{lstlisting}
Preciso de um esqueleto de API backend.
Requisitos:
1. Servidor Express, porta 3001
2. Banco de dados SQLite, tabela records
3. Campos da tabela records:
   id (TEXT PRIMARY KEY)
   amount (REAL)
   type (TEXT: 'income'|'expense')
   category (TEXT)
   note (TEXT)
   date (TEXT: YYYY-MM-DD)
4. Endpoints CRUD:
   GET /api/records - obter todos os registros
   POST /api/records - adicionar registro
   PUT /api/records/:id - atualizar registro
   DELETE /api/records/:id - excluir registro
5. Middleware de tratamento de erros
Dê a estrutura completa do código backend.
\end{lstlisting}

Repare nas características dessa instrução: ela descreve apenas o passo atual (o esqueleto do back-end), sem misturar requisitos de front-end; lista explicitamente os campos e os endpoints (dando restrições suficientes à IA), mas deixa a implementação concreta a cargo da IA; e segue o princípio Goal-Driven Execution — dizer 'o que fazer', e não 'como fazer'.

Assim que a IA entregar o código, valide na hora: rode o servidor e teste cada endpoint com curl ou Postman para ver se funciona. Confirmado que está tudo certo, abra uma nova conversa para o Passo 2 (páginas básicas do front-end). Cada passo segue o ciclo 'instrução – validação – nova conversa'.

\booksection{Projeto 3: aplicação turbinada com IA (nível avançado)}

\booksubsection{Concepção do projeto}

Objetivo: uma aplicação que consome a API de grandes modelos — por exemplo, um 'assistente de estudos com IA', em que o usuário insere o conteúdo de estudo e a IA gera resumos, conecta tópicos relacionados e recomenda materiais. O diferencial deste projeto: a sua aplicação, ela mesma, chama a IA, e você vai precisar lidar com questões novas como chamadas de API, engenharia de prompt e tratamento de erros.

A lição mais importante deste projeto é: quando você 'usa a IA para escrever o código que chama a IA', precisa dar atenção redobrada à engenharia de prompt. Os prompts que a sua aplicação usa ao chamar a API do grande modelo também devem seguir os princípios discutidos neste guia — clareza, concisão e orientação a objetivos.

\booksubsection{Pontos técnicos decisivos}

Este projeto envolve alguns pontos técnicos que você não encontrou nos dois projetos anteriores:

\begin{itemize}
  \item Gerenciamento da chave da API: a chave jamais deve ficar escrita no código; gerencie-a com variáveis de ambiente;
  \item Tratamento de respostas em streaming: as respostas dos grandes modelos normalmente chegam em streaming e precisam ser exibidas em tempo real, trecho a trecho;
  \item Controle do consumo de Tokens: os grandes modelos cobram por token, então o consumo precisa ser controlado no nível da aplicação;
  \item Gerenciamento de templates de prompt: funcionalidades diferentes usam templates de prompt diferentes, o que exige uma gestão estruturada.
\end{itemize}

Todos esses pontos técnicos podem ser implementados com a ajuda da IA, mas você precisa entender o 'porquê' de cada um — por que a chave não pode ficar no código? Por que é preciso streaming? Por que o consumo de tokens precisa ser controlado? Só entendendo o 'porquê' você consegue julgar se a solução proposta pela IA faz sentido.

\booksection{Checklist universal do fluxo prático}

Não importa o tamanho do projeto: a seguir está um checklist universal do fluxo prático de programação com IA. Sempre que começar um projeto novo, percorra essa lista para garantir que nenhum passo essencial fique de fora.

\begin{center}
\begin{tabularx}{\linewidth}{|X|X|X|}
\hline
\textbf{Etapa} & \textbf{O que fazer} & \textbf{Princípio correspondente} \tabularnewline
\hline
1. Pensar nos requisitos & Percorrer os cinco elementos de requisitos e esclarecer o que será feito & Think Before Coding \tabularnewline
2. Decompor os requisitos & Dividir a demanda grande em pequenos passos que possam ser validados de forma independente & Simplicity First \tabularnewline
3. Escrever o PRD & Organizar em um documento estruturado para servir de referência & Goal-Driven \tabularnewline
4. Configurar o arquivo de memória & Criar o CLAUDE.md ou arquivo equivalente & Harness Engineering \tabularnewline
5. Implementar passo a passo & Uma conversa por passo: instrução – validação – próximo passo & Surgical Changes \tabularnewline
6. Testar e validar & Testar a correção da funcionalidade imediatamente após cada passo & Iterate \& Verify \tabularnewline
7. Revisar o código & Verificar sanidade da arquitetura, segurança e facilidade de manutenção & Iterate \& Verify \tabularnewline
8. Iterar e otimizar & Melhorar em iterações conforme os resultados dos testes e da revisão & Iterate \& Verify \tabularnewline
\hline
\end{tabularx}
\end{center}

\begin{center}\small\color{gray}Checklist universal do fluxo prático de programação com IA\end{center}

\begin{tcbpractice}{Mãos à obra: conclua o 'Projeto 1' nesta semana}
Trate o 'site de ferramentas pessoais' da Seção 1 do Capítulo 44 como a sua primeira tarefa de programação com IA: dedique 5 minutos aos cinco elementos, envie a instrução à IA (qualquer ferramenta serve) e faça a aceitação com as três camadas de verificação. Termine em até 30 minutos — e a sua primeira obra de programação com IA estará pronta. Guarde bem o resultado: é a lembrança da sua passagem de 'leitor' a 'praticante'.

\end{tcbpractice}

\begin{tcbsummary}{Resumo do capítulo}
\begin{itemize}
  \item Projeto 1 (iniciante): site de ferramentas só com front-end — vivencie o fluxo completo; dá para terminar em 30 minutos
  \item Projeto 2 (intermediário): sistema de controle financeiro full stack — PRD + arquivo de memória + implementação em 5 passos, com algo 'que funciona' ao fim de cada um
  \item Projeto 3 (avançado): assistente de estudos com IA — use a IA para escrever o código que chama a IA, com atenção à engenharia de prompt
  \item Fluxo universal em oito passos: requisitos decomposição PRD arquivo de memória implementação passo a passo testes revisão iteração
  \item Cada passo segue o ciclo 'instrução – validação – nova conversa'
\end{itemize}

\end{tcbsummary}
